
 To relax the constraint $x_1+x_2=2$, we first write it as $2-x_1-x_2
 = 0$. Indeed, the quantity $2-x_1-x_2$ quantifies the amount of
 violation of the constraint. If it is zero, the constraint is not
 violated. If it is different from zero, it is.  We introduce a penalty that is proportional to the violation of the constraint, with a proportionality factor $\lambda$. The new optimization problem is

 \[
 \min_{x \in \mathbb{R}^2} x_1+4x_2 + \lambda(2-x_1-x_2)
 \]
 subject to 
\begin{align*}
x_1 & \geq 0,\\
x_2 & \geq 0.
\end{align*}

 We examine the different cases corresponding to the different values of $\lambda$.

\begin{itemize}
\item If $\lambda=0$, the problem is
 \[
 \min_{x \in \mathbb{R}^2} x_1+4x_2
 \]
 subject to 
\begin{align*}
x_1 & \geq 0,\\
x_2 & \geq 0.
\end{align*}
Its optimal solution is  $x^*=(0,0)^T$ with a value of $0$. This
solution violates the constraint of the original problem, and the
optimal value is lower. 

\item If $\lambda=1$, the problem is
 \[
 \min_{x \in \mathbb{R}^2} 2+3x_2 
 \]
 subject to 
\begin{align*}
x_1 & \geq 0,\\
x_2 & \geq 0.
\end{align*}
Any feasible $x$ such that $x_2=0$ is an optimal solution to the
problem with a value of $2$, which is the same optimal value as the
original problem, with the constraint. Because of the penalty,
there is no incentive to violate the constraint, as it does not allow
to improve the solution of the original problem. 

\item If $\lambda=2$, the problem is
 \[
 \min_{x \in \mathbb{R}^2} 4-x_1+ 2x_2
 \]
 subject to 
\begin{align*}
x_1 & \geq 0,\\
x_2 & \geq 0.
\end{align*}
The problem is not bounded, as the value of the objective function
goes to $-\infty$ when $x_1$ goes to $+\infty$.The value of the
penalty parameter is not appropriate, as it transforms a well-defined
problem into an unbounded problem. 

\end{itemize}

